\section{Introduction}
\label{sec:intro}

Captchas have evolved from distorted text images to dynamic, multi-step challenges
designed to resist automation while remaining trivially solvable by humans.
A recent class of video-editing captchas presents users with a short clip and a natural-language
instruction such as ``reverse the video, then flip it horizontally.''
Following these instructions reveals an embedded alphanumeric key that serves as the captcha answer.
A human solves this in seconds; an automated agent must parse the instruction,
execute the correct sequence of video transformations, and extract the now-visible key---
a compound task spanning language understanding, procedural execution, and visual perception.

Despite the proliferation of video-capable VLMs~\citep{gpt4o,gemini,qwen2vl},
no benchmark systematically tests this compound capability.
Existing video benchmarks target temporal reasoning~\citep{videomme,longvideobench}
or egocentric understanding~\citep{ego4d}, not instruction-conditioned video manipulation.
Video-editing benchmarks~\citep{vbench,editbench} evaluate generative models rather than agents.
CIPHER fills this gap.

\paragraph{Contributions.}
\begin{enumerate}
  \item \textbf{Task formulation.} We define the video-editing captcha solving task
        as a three-stage pipeline: instruction parsing (S1), ffmpeg execution (S2),
        and key extraction (S3), enabling fine-grained per-stage diagnostics.

  \item \textbf{Benchmark.} CIPHER provides 250 procedurally generated samples
        across five difficulty levels (L1--L5), varying key visibility,
        operation count (1--4 ops), and instruction ambiguity (exact to compositional).
        Backgrounds are drawn from Kinetics-400 to ensure visual diversity.

  \item \textbf{Metrics.} We propose \emph{Pipeline-F1}, a weighted combination of
        per-stage scores ($0.25 \cdot \text{S1} + 0.25 \cdot \text{S2} + 0.50 \cdot \text{S3}$),
        alongside strict exact match (EM) and character-level F1.

  \item \textbf{Baselines.} We benchmark OCR-only, Claude Haiku, GPT-4o-mini,
        Gemini Flash, and local open-weight models (Qwen2-VL-7B/72B),
        revealing that S1 parsing saturates quickly while S3 extraction
        on natural backgrounds remains the binding bottleneck.
\end{enumerate}
